%!TEX TS-program = xelatex
\documentclass[]{friggeri-cv}
\usepackage{afterpage}
\usepackage{hyperref}
\usepackage{color}
\usepackage{xcolor}
\usepackage{smartdiagram}
\usepackage{fontspec}
% if you want to add fontawesome package
% you need to compile the tex file with LuaLaTeX
% References:
%   http://texdoc.net/texmf-dist/doc/latex/fontawesome/fontawesome.pdf
%   https://www.ctan.org/tex-archive/fonts/fontawesome?lang=en
%\usepackage{fontawesome}
\usepackage{metalogo}
\usepackage{dtklogos}
\usepackage[utf8]{inputenc}
\usepackage{tikz}
\usetikzlibrary{mindmap,shadows}
\hypersetup{
    pdftitle={},
    pdfauthor={},
    pdfsubject={},
    pdfkeywords={},
    colorlinks=false,           % no lik border color
    allbordercolors=white       % white border color for all
}
\smartdiagramset{
    bubble center node font = \footnotesize,
    bubble node font = \footnotesize,
    % specifies the minimum size of the bubble center node
    bubble center node size = 0.5cm,
    %  specifies the minimum size of the bubbles
    bubble node size = 0.5cm,
    % specifies which is the distance among the bubble center node and the other bubbles
    distance center/other bubbles = 0.3cm,
    % sets the distance from the text to the border of the bubble center node
    distance text center bubble = 0.5cm,
    % set center bubble color
    bubble center node color = pblue,
    % define the list of colors usable in the diagram
    set color list = {lightgray, materialcyan, orange, green, materialorange, materialteal, materialamber, materialindigo, materialgreen, materiallime},
    % sets the opacity at which the bubbles are shown
    bubble fill opacity = 0.6,
    % sets the opacity at which the bubble text is shown
    bubble text opacity = 0.5,
}

\addbibresource{bibliography.bib}
\RequirePackage{xcolor}
\definecolor{pblue}{HTML}{0395DE}

\begin{document}
\header{Amarjeet Singh }{Kapoor}
      {Computer Engineer}
      
% Fake text to add separator      
\fcolorbox{white}{gray}{\parbox{\dimexpr\textwidth-2\fboxsep-2\fboxrule}{%
.....
}}

% In the aside, each new line forces a line break
\begin{aside}
\includegraphics[scale=0.30]{amar.png}
      \section{Mail}
    \href{mailto:amarjeet.kapoor1@gmail.com}{\textbf{amarjeet.kapoor1@}\\gmail.com}
  \section{Tel}
    +91 8568988521
  \section{Git and Web}
    \href{http://github.com/amarjeetkapoor1}{\textbf{github.com/amarjeet \\ kapoor1}}
    \href{http://amarjeetkapoor1.wordpress.com}{\textbf{amarjeetkapoor1.\\wordpress.com}}
  \section{Programming \\ Language}
    \smartdiagram[bubble diagram]{
        \textbf{C/C++},
        \textbf{\hspace{1.5mm} Python \hspace{1.5mm} \vspace{3mm}},
        \textbf{TeX},
        \textbf{HTML}\\\textbf{JS},
        \textbf{Bash},
        \textbf{SQL},
        \textbf{Octave},
        Mark\\down
    }
    
  \section{Personal Skills}
    \smartdiagram[bubble diagram]{
        \textbf{Team}\\\textbf{Player},
        \textbf{Curiosity},
        \textbf{Problem}\\\textbf{Solving},
        \textbf{Transcend}\\ {Lang.},
        \textbf{Logical}\\ {Thinking},
        \textbf{Self}\\ {Learner},
        \textbf{Organize}
    }
    ~
  \section{Technologies \\ Used}
        \textbf{OpenCV}
        \textbf{LaTeX}
        \textbf{Django}
        \textbf{QT}
        \textbf{SageMath}
        \textbf{Doxygen}
        \textbf{Flex}
        \textbf{Bison}
        \textbf{GIT}
        \textbf{GRASS GIS} 
     ~      
\end{aside}

\section{Objective}
To make my story matter to world.
\section{Education}
\begin{entrylist}
  \entry
    {2013 - 2017}
    {Bachelor's Degree in Computer Science and Engineering}
    {78 \%}
    {Guru Nanak Dev Engineering College}
  \entry
    {2013}
    {Higher Secondary Examination}
    {90.8 \%}
    {Green Land Sr. Sec. Public School (CBSE)}
    \entry
    {2011}
    {Matriculation}
    {8.6 CGPA}
    {Green Land Sr. Sec. Public School (CBSE) }
\end{entrylist}

\section{Projects}
\begin{entrylist}
  \entry
    {}
    {OpenSCAD Customizer}
    {OpenSCAD}
    {
        This is a Google Summer of Code (GSoC) project under OpenSCAD community. It intends to provide syntax for creating a User Interface for modifying 3D models.
        \textit{\href{https://github.com/openscad/openscad/tree/gsoc2016}{https://github.com/opencad/openscad/tree/gsoc2016}}
    }
  \entry
    {}
    {CivilOctave}
    {6-Week Training}
    {
        It is a project for studying the Dynamic of Structure and provide users with
results from which decision related to stability of structure can be taken.  
        \textit{\href{https://github.com/amarjeetkapoor1/CivilOctave}{https://github.com/amarjeetkapoor1/CivilOctave}}
    }
    \entry
    {}
    {Sim (Structure information modeling)}
    {TCC}
    {
       This is a large undertaking for which work is still being done. This software
intends to provide users with ability to store information of a structure and
provide full analysis of that structure.
        \textit{\href{https://github.com/amarjeetkapoor1/Sim}{https://github.com/amarjeetkapoor1/Sim}}
    }
    \entry
    {}
    {Image processing}
    {}
    {
    Programs to Finds the number of significant colors and to find the objects of
particular shape and of particular color from camera.
        \textit{\href{https://github.com/amarjeetkapoor1/image-process}{https://github.com/amarjeetkapoor1/image-process}}
    }
    
    \entry
    {}
    {TheRoadProject}
    {TCC}
    {
    This is the automation of process of finding the sections of road at particular
distance and then calculating the RL values.
    \textit{\href{https://github.com/amarjeetkapoor1/TheRoadProject}{ https://github.com/amarjeetkapoor1/TheRoadProject}
    }}
    
        \entry
    {}
    {AppRoxEr}
    {DeeCoders}
    {
     AppRoxEr is an android app based on HTML parsing that serves the student of GNDEC for the purposes of academic attendance, time table, etc.
    }
    
\end{entrylist}


\section{Achievements}
\begin{itemize}
\item Google Summer of Code, 2016.
\item Successfully collaborating with open-source projects. 
\item Won District level medals in table tennis.
\end{itemize}

\end{document}
